% =============================================================================
% Глава 6. Адверсариально-устойчивая защита нейросетевых систем
% навигации БПЛА на основе динамических графовых нейронных сетей
% -----------------------------------------------------------------------------
% Файл предназначен для включения в PhD_Thesis_Integrated.tex через \include.
% Помимо пакетов, уже подключённых в основном файле, для корректной
% компиляции главы на русском языке требуется добавить в преамбулу
% диссертации:
%
%     \usepackage[T2A]{fontenc}                 % кириллица
%     \usepackage[english,russian]{babel}       % основной язык — русский
%     \usepackage{pgfplots}\pgfplotsset{compat=1.18}
%     \usetikzlibrary{positioning,arrows.meta,shapes.geometric,
%                     calc,fit,backgrounds,patterns}
%
% Все используемые в главе команды (\R, \E, \Loss, \argmin, \argmax,
% theorem-окружения) либо стандартны, либо уже определены в основном файле.
% Дополнительно вводятся локальные команды:
\providecommand{\Adv}{\mathcal{A}}
\providecommand{\Def}{\mathcal{D}}
\providecommand{\Loss}{\mathcal{L}}
\providecommand{\eps}{\varepsilon}
\providecommand{\Norm}[1]{\left\lVert#1\right\rVert}
\providecommand{\sign}{\operatorname{sign}}
% =============================================================================

\chapter{Адверсариально-устойчивая защита нейросетевых систем навигации
беспилотных летательных аппаратов на основе динамических графовых
нейронных сетей}
\label{ch:uav_defense}

\section{Введение к главе}
\label{sec:uav_intro}

Развитие предложенной в главах~4 и~5 методологии динамических графовых
нейронных сетей (DG-GNN) и оптимально-транспортной федеративной защиты
позволило получить устойчивые количественные результаты на платформе
\textsf{RobustIDPS.ai}~\cite{robustidps_zenodo} на корпусе из 84{,}2~млн
размеченных сетевых событий. В то же время постановочный документ
проекта~\cite{uav_proposal_2026} ставит задачу переноса доказанного
аппарата в принципиально иную предметную область~--- задачу обеспечения
адверсариальной устойчивости нейросетевых компонентов системы навигации
беспилотных летательных аппаратов (БПЛА) в условиях радиоэлектронной
борьбы (РЭБ) и активного состязательного машинного обучения.

Актуальность задачи подтверждается тремя независимыми группами
наблюдений. Во-первых, материалы открытых военных
оценок~\cite{wapo_excalibur_2024} фиксируют падение доли поражений
управляемого артиллерийского снаряда M982~Excalibur с уровня выше
50\,\% до менее 10\,\% в течение нескольких месяцев применения в зоне
действия российских средств подавления глобальных навигационных
спутниковых систем (ГНСС). Во-вторых, по данным
\cite{inside_uas_2025}, до 31\,\% боевых вылетов украинских FPV-БПЛА
теряется вследствие радиоэлектронного подавления, а Латвийское
управление электронных коммуникаций зафиксировало в 2024~году
820~инцидентов воздействия на ГНСС против 26 в~2022~году. В-третьих,
академическая литература подтверждает, что прицельные адверсариальные
заплатки против аэрофотодетекторов снижают средние значения метрики
average precision до~87{,}86~процентных пунктов~\cite{lian_2022}, а
атаки методом Basic Iterative Method на политики глубокого обучения с
подкреплением для пилотирования БПЛА уменьшают долю успешных проходов
полосы препятствий с 97\,\% до 35\,\%~\cite{hickling_2023}.

С точки зрения нормативной базы, в России Постановление Правительства
№\,1701 от 30~ноября~2024~года~\cite{decree_1701}, вступающее в силу
поэтапно с~1~марта~2025 по~2027~годы, обязывает оснащать БПЛА массой
свыше~250\,г сертифицированными средствами связи, удалённой
идентификации и криптографической защиты информации. Действующий
\mbox{ГОСТ~Р~72118-2025} регламентирует формализацию моделей угроз для
систем искусственного интеллекта; международный контур формируют
\mbox{NIST~AI~RMF~1.0}~\cite{nist_rmf}, акт ЕС об~ИИ и
\mbox{DO-326A/ED-202A}~\cite{do_326a}. Все они требуют
\emph{сертифицируемой, аудируемой и приватной} робастности, то есть
именно тех свойств, которые предоставлены диссертационным аппаратом на
реляционных данных в~главах~4--5.

Цель главы состоит в распространении валидированных в основной части
работы методов \textsf{M1}--\textsf{M7} и фундаментальной модели
\textsf{CyberSecLLM} на предметную область навигации БПЛА с сохранением
доказанных сертификатов устойчивости и приватности. Сформулированы
следующие задачи: (1)~формализовать операционный граф БПЛА-роя как
непрерывно-временной динамический объект; (2)~получить математический
аппарат шести семейств атак, регламентированных эталонным набором
проекта~\cite{uav_proposal_2026}, в едином графовом представлении;
(3)~отобразить методы основной части на конкретные механизмы защиты в
трёхуровневой архитектуре \emph{борт~--- дронопорт~--- облако};
(4)~получить априорные сертификаты Липшица--Гронуолла и рандомизованного
сглаживания; (5)~эмпирически проверить работоспособность всего конвейера
на референсной реализации Phase~A.

\subsection*{Научная новизна главы}
Научная новизна главы состоит в следующем.
\begin{enumerate}\setlength\itemsep{0.1em}
\item Впервые предложен единый \emph{граф-нейросетевой} конвейер защиты
БПЛА, объединяющий непрерывно-временную динамику узлов
(\textsf{CT-TGNN}), селективное пространственно-состояний моделирование
(\textsf{MambaShield}), визант-устойчивую федеративную агрегацию
(\textsf{FedLLM-API}) и игровую стохастическую сертификацию
(\textsf{FedGTD}).
\item Получены количественные сертификаты робастности
(теоремы~\ref{thm:gronwall_uav}--\ref{thm:mwu_uav}), переносящие
доказанные в~главах~4--5 гарантии на предметную область навигации БПЛА.
\item Построена матрица соответствий $3\times4$
\emph{условие~--- требование} и доказана покрываемость всех её
двенадцати клеток как минимум одним сертифицированным методом
основной части работы.
\item Подготовлена референсная реализация фазы~A
(см.~\S\ref{sec:uav_phase_a}), документирующая практическую
работоспособность конвейера и интеграцию с производственной платформой
\textsf{RobustIDPS.ai}.
\end{enumerate}

\subsection*{Структура главы}
В~\S\ref{sec:uav_problem} вводится формальная постановка задачи и
графовая модель оперативной обстановки БПЛА. В~\S\ref{sec:uav_threat}
описываются шесть семейств атак, регламентированных эталонным набором
проекта. В~\S\ref{sec:uav_math} приводятся четыре сертификата
устойчивости. В~\S\ref{sec:uav_arch} представлена трёхуровневая
архитектура унифицированного защитного фреймворка и матрица
\emph{условие~--- требование}. В~\S\ref{sec:uav_mapping} отображены
методы \textsf{M1}--\textsf{M7} на конкретные механизмы защиты в каждой
из трёх прикладных областей UAV-стека. В~\S\ref{sec:uav_algorithm}
формализуется унифицированная процедура обучения.
\S\ref{sec:uav_phase_a} содержит результаты экспериментальной апробации
фазы~A и сопоставление с показателями, полученными ранее на платформе
\textsf{RobustIDPS.ai}. В~\S\ref{sec:uav_regulation} проанализировано
соответствие предлагаемого решения нормативной базе. Выводы по главе
сформулированы в~\S\ref{sec:uav_conclusion}.

% =============================================================================
\section{Постановка задачи и графовая модель оперативной обстановки}
\label{sec:uav_problem}

\subsection{Операционный граф БПЛА}

Пусть в момент времени $t$ группа из $n$ БПЛА совместно с одной или
несколькими наземными станциями управления и дронопортами образует
\emph{операционный граф}
\begin{equation}
\mathcal{G}_t \;=\; \bigl(V_t,\,E_t,\,X_t,\,A_t\bigr),
\label{eq:uav_graph}
\end{equation}
в котором $V_t$~--- множество узлов, объединяющее воздушные БПЛА
$u_i$, дронопорты $p_j$ и наземные конечные точки управления $g_k$;
$E_t$~--- направленные рёбра, кодирующие активные радио-, оптические,
инерциальные и ГНСС-каналы; $X_t \in \R^{n\times d}$~--- матрица
признаков узлов, агрегирующая телеметрию (состояние двигателей,
целостность прошивки, состояние ИНС/ГНСС, состояние полезной нагрузки);
$A_t \in \R^{n\times n}$~--- весовая матрица смежности, изменяющаяся во
времени вследствие активации, затухания и подавления каналов связи.

Защитник развёртывает на каждом узле нейросетевую функцию предсказания
\begin{equation}
f_{\theta}\!\colon\,\bigl(X_{\le t},\,A_{\le t}\bigr) \,\longrightarrow\,
\hat{y}_t,
\end{equation}
выходом которой служит решение задачи восприятия (обнаружение/слежение),
навигации (опорная точка, угловое положение) или состояния миссии.
Атакующий $\Adv$ выбирает возмущение $\delta$ из ограничивающего
множества $\mathcal{S}$, отвечающего одному из шести семейств эталонного
набора~\cite{uav_proposal_2026}; формально задача атакующего имеет вид
\begin{equation}
\delta^{*} \;=\; \argmax_{\delta\in\mathcal{S}}
\Loss\!\bigl(f_{\theta}(X_{\le t}+\delta,\,A_{\le t}),\,y_t\bigr).
\label{eq:uav_advmax}
\end{equation}
Двойственная задача защитника состоит в построении такой функции
$f_{\theta}$, для которой при любом $\delta\in\mathcal{S}$ выполняется
\emph{количественная} оценка чувствительности
\begin{equation}
\Norm{f_{\theta}(X+\delta) - f_{\theta}(X)} \;\le\; \rho(\eps),
\label{eq:uav_cert}
\end{equation}
где $\rho$~--- известная функция радиуса возмущения~$\eps$.

\subsection{Отображение классов данных проектной модели на граф}
Данные требования проекта~(п.~4.2.1)~\cite{uav_proposal_2026}
организуются в графе следующим образом. \emph{Полётная миссия} образует
целевую разметку в узлах (плановые путевые точки, режимы и временные
окна). \emph{Состояние БПЛА} формирует признаки узлов $X_t$ с
криптографическим заверением целостности по \mbox{ГОСТ~Р~34.10-2012}.
\emph{Геоданные} входят в веса рёбер (расстояние до запретных зон) и в
структурные ограничения матрицы $A_t$. \emph{Данные наземной
инфраструктуры} связываются с узлами дронопортов. \emph{Телеметрия}
поступает в качестве входов с временными отметками в
непрерывно-временную динамику. \emph{Метаданные} обеспечивают
поставщиков подлинности цифровых подписей для последующего аудита.

% =============================================================================
\section{Эталонный набор атак}
\label{sec:uav_threat}

В соответствии с п.~4.4.1.1 постановочного
документа~\cite{uav_proposal_2026} рассматриваются три класса атак.

\paragraph{Атаки белого ящика.}
Атака \emph{FGSM}~\cite{goodfellow_fgsm} в едином направлении вдоль
знака градиента
\begin{equation}
X_{\text{adv}} \;=\; X + \eps\,\sign\!\bigl(\nabla_{X}\Loss(\theta,X,y)\bigr).
\label{eq:uav_fgsm}
\end{equation}
Многошаговая итеративная атака \emph{PGD}~\cite{madry_pgd} с проекцией
на $\ell_p$-шар
\begin{equation}
X_{t+1} \;=\; \Pi_{X+\mathcal{S}}\!\Bigl(X_t + \alpha\,\sign\bigl(\nabla_X\Loss\bigr)\Bigr).
\label{eq:uav_pgd}
\end{equation}
Атака Карлини--Уэйнера \emph{CW}~\cite{carlini_cw} минимизирует норму
возмущения при ограничении на разрыв логитов:
\begin{equation}
\min_{\delta}\ \Norm{\delta}_p + c\cdot f_{6}(X+\delta),\qquad
f_{6}(X') \;=\; \max\!\bigl(\!\max_{i\neq t}\!Z(X')_i - Z(X')_t,\,-\kappa\bigr).
\label{eq:uav_cw}
\end{equation}
Условие принадлежности $X+\delta$ единичному гиперкубу обеспечивается
переменой переменных $\delta = \tfrac{1}{2}(\tanh w + 1) - X$.

\paragraph{Атаки чёрного ящика.}
Атака HopSkipJump~\cite{chen_hsj} использует оценку градиента методом
Монте-Карло
\begin{equation}
\widetilde{\nabla}S(X_t,\delta_t) \;=\; \frac{1}{B}\sum_{b=1}^{B}
\phi_{X^{*}}\!\bigl(X_t+\delta_t u_b\bigr)\,u_b,\qquad
u_b\sim\mathrm{Unif}(\mathcal{S}^{d-1}),
\label{eq:uav_hsj}
\end{equation}
где $\phi$~--- решающий оракул. Атака BoundaryAttack~\cite{brendel_ba}
реализует случайное блуждание вдоль границы решения.

\paragraph{Атака на этапе обучения~--- отравление данных.}
Двухуровневая задача
\begin{equation}
\max_{D'_{c}}\ \Loss(D_{\text{val}},\theta^{\!*})\quad
\text{при}\quad
\theta^{\!*}\!\in\!\argmin_{\theta}\Loss(D_{\text{tr}}\cup D'_{c},\theta).
\label{eq:uav_poison}
\end{equation}
В качестве канонического примера используется атака
``Poison~Frogs!''~\cite{shafahi_poison} с чистыми метками и точкой
коллизии в пространстве признаков.

% =============================================================================
\section{Сертификаты устойчивости}
\label{sec:uav_math}

Перенесём сертификаты, полученные в~главах~4--5, на построенный
графовый объект~\eqref{eq:uav_graph}. Все четыре утверждения настоящей
секции переиспользуют доказательства основной части без изменений; для
полноты главы приведены ключевые конструкции.

\begin{theorem}[Сертификат Липшица--Гронуолла для \textsf{CT-TGNN}]
\label{thm:gronwall_uav}
Пусть динамика скрытого состояния узлов на непрерывно-временном
операционном графе БПЛА удовлетворяет уравнению
\begin{equation}
\frac{\mathrm{d}h_v(t)}{\mathrm{d}t} \;=\;
f_{\theta}\!\bigl(h_v(t),A(t),X(t),t\bigr),
\label{eq:uav_dyn}
\end{equation}
причём оператор $f_{\theta}$ удовлетворяет условию Липшица с константой
$L_g$ по первому аргументу. Тогда для любых двух входов $x_1,x_2$
выполняется неравенство
\begin{equation}
\Norm{h_1(T)-h_2(T)} \;\le\; e^{L_g T}\,\Norm{x_1-x_2}.
\label{eq:uav_gronwall}
\end{equation}
\end{theorem}
\begin{proof}[Эскиз доказательства]
Стандартное следствие из неравенства Гронуолла, применённого к функции
$z(t) = \Norm{h_1(t)-h_2(t)}$ с использованием оценки $\dot{z}\le L_g z$
и начального условия $z(0) = \Norm{x_1-x_2}$.\qedhere
\end{proof}

\begin{theorem}[Сертифицированный $\ell_2$-радиус через
рандомизованное сглаживание~\cite{cohen_smoothing}]
\label{thm:rs_uav}
Пусть $g(x) = \argmax_{c}\Pr_{\eta\sim\mathcal{N}(0,\sigma^{2}I)}\!
\bigl[f(x+\eta) = c\bigr]$, и пусть для некоторого класса $c_A$
\[
\Pr\!\bigl[f(x+\eta) = c_A\bigr] \ge p_A,\qquad
\max_{c\ne c_A}\Pr\!\bigl[f(x+\eta) = c\bigr] \le p_B.
\]
Тогда для любого возмущения $\delta$ с $\Norm{\delta}_2 \le R$
\[
g(x+\delta) \;=\; c_A,\qquad
R \;=\; \tfrac{\sigma}{2}\!\bigl(\Phi^{-1}(p_A)-\Phi^{-1}(p_B)\bigr).
\]
\end{theorem}

\begin{theorem}[PAC-байесовская оценка обобщения для
\textsf{MambaShield}]\label{thm:pac_uav}
Для произвольного априорного распределения~$\pi$, апостериорного~$\rho$
и любого $\delta\in(0,1)$ с вероятностью не менее $1-\delta$ по
выборке независимых одинаково распределённых наблюдений размера~$N$
выполняется
\begin{equation}
R(\theta) \;\le\; \hat R(\theta) +
\sqrt{\frac{\mathrm{KL}(\rho\,\|\,\pi)+\ln(2N/\delta)}{2N}}.
\label{eq:uav_pacb}
\end{equation}
\end{theorem}

\begin{theorem}[Регретный сертификат мультипликативных весов
для \textsf{FedGTD}]\label{thm:mwu_uav}
Для защитника, выбирающего стратегию $i\in\{1,\ldots,|S|\}$ с весами
$w_i^{(t+1)} = w_i^{(t)}\exp\!\bigl(-\eta\,\ell_i^{(t)}\bigr)$ и шагом
$\eta = \sqrt{\ln|S|/T}$, накопленный регрет ограничен
\begin{equation}
R(T) \;\le\; \sqrt{T\ln|S|}.
\label{eq:uav_mwu}
\end{equation}
\end{theorem}

Теоремы~\ref{thm:gronwall_uav}--\ref{thm:mwu_uav} образуют замкнутую
систему гарантий: \eqref{eq:uav_gronwall} ограничивает чувствительность
выхода при произвольном входном возмущении из шара радиуса~$\eps$;
\eqref{eq:uav_pacb} переносит эту оценку с обучающей на тестовую
выборку; рандомизованное сглаживание выдаёт сертифицированный радиус
$R$, не зависящий от внутренней структуры классификатора; наконец,
теорема~\ref{thm:mwu_uav} гарантирует сублинейный регрет защитника при
любой адаптивной стратегии атакующего.

% =============================================================================
\section{Архитектура унифицированного защитного фреймворка}
\label{sec:uav_arch}

Конструктивно фреймворк реализован в виде трёхуровневой архитектуры
\emph{борт~--- дронопорт~--- облако} (см.\ рис.~\ref{fig:uav_arch}),
наследующей программный контур платформы~\textsf{RobustIDPS.ai} и
расширяющей его модально-агностической библиотекой
\textsf{robustnn-core} с интерфейсом \texttt{UAVGraphBatch}. На бортовом
уровне выполняются методы \textsf{M1~CT-TGNN}, \textsf{M4~MambaShield} и
\textsf{M6~UC-HGP}, обеспечивающие низкоресурсный непрерывно-временной
вывод. Уровень дронопорта реализует методы \textsf{M2~FedLLM-API} (с
визант-устойчивой агрегацией) и \textsf{M7~FedGTD} (стохастическая игра
защитника против адаптивного источника помех). Облачный уровень
ответствен за зашитный аудит миссий \textsf{CyberSecLLM} и многоуровневую
ретроспективу \textsf{M3~TripleE-TGNN}.

\begin{figure}[t]
\centering
\begin{tikzpicture}[font=\footnotesize,
  tier/.style={rectangle,rounded corners=2pt,draw=primaryblue!80,
               thick,fill=primaryblue!8,minimum height=1.1cm,
               minimum width=2.6cm,align=center,inner sep=3pt,
               font=\small\bfseries},
  mod/.style={rectangle,rounded corners=1.5pt,draw=darkgreen!80,
              thick,fill=darkgreen!10,minimum height=0.45cm,
              minimum width=2.4cm,inner sep=2pt,font=\footnotesize\bfseries},
  threat/.style={rectangle,rounded corners=1.5pt,draw=accentorange!80,
                 thick,fill=accentorange!15,minimum height=0.45cm,
                 minimum width=2.6cm,inner sep=2pt,font=\footnotesize},
  link/.style={->,>=Stealth,thick,primaryblue!75},
  attk/.style={->,>=Stealth,thick,accentorange,dashed}]
\node[tier]  (edge)  at (0,0)    {Борт БПЛА\\($<\!5\,$Вт SoC)};
\node[tier]  (port)  at (4.0,0)  {Дронопорт\\(GPU-узел)};
\node[tier]  (cloud) at (8.0,0)  {Облако / SOC\\(глобальный)};
\node[mod, above=3pt of edge]  (m1)  {M1 CT-TGNN};
\node[mod, above=3pt of m1]    (m4)  {M4 MambaShield};
\node[mod, above=3pt of m4]    (m6)  {M6 UC-HGP};
\node[mod, above=3pt of port]  (m2)  {M2 FedLLM-API};
\node[mod, above=3pt of m2]    (m7)  {M7 FedGTD};
\node[mod, above=3pt of m7]    (m5)  {M5 S-Trans.};
\node[mod, above=3pt of cloud] (m3)  {M3 TripleE-TGNN};
\node[mod, above=3pt of m3]    (cll) {CyberSecLLM};
\node[mod, above=3pt of cll]   (rep) {RobustIDPS.ai};
\node[threat, below=6pt of edge]  (t1) {FGSM / PGD / CW};
\node[threat, below=6pt of port]  (t2) {Византийцы / Отравл.};
\node[threat, below=6pt of cloud] (t3) {HopSkip / Boundary};
\draw[link] (edge.east) -- (port.west);
\draw[link] (port.east) -- (cloud.west);
\draw[attk] (t1) -- (edge);
\draw[attk] (t2) -- (port);
\draw[attk] (t3) -- (cloud);
\node[font=\footnotesize\itshape,accentorange!90]
  at ($(t2.south)+(0,-0.32)$)
  {Эталонный набор атак (п.~4.4.1.1)};
\end{tikzpicture}
\caption[Трёхуровневая архитектура защитного фреймворка.]%
{Трёхуровневая архитектура унифицированного защитного фреймворка
\emph{борт~--- дронопорт~--- облако}. На каждом уровне развёрнут
подмножество методов \textsf{M1}--\textsf{M7} и фундаментальная
модель \textsf{CyberSecLLM}; атаки эталонного набора (п.~4.4.1.1)
действуют преимущественно на нижнем уровне борта, тогда как
визант-устойчивая агрегация и облачный аудит сосредоточены на
вышестоящих уровнях.}
\label{fig:uav_arch}
\end{figure}

Центральным организующим объектом служит \emph{матрица соответствия}
(табл.~\ref{tab:uav_matrix}): три эксплуатационных условия
(адверсариальные, нестационарные, распределённые) против четырёх
требований (робастность, точность, эффективность, приватность). Каждой
клетке поставлен в соответствие как минимум один метод
основной части, выдающий количественный сертификат.

\begin{table}[t]
\centering
\caption{Матрица соответствия \emph{условие~--- требование} с
указанием первичных методов для предметной области БПЛА.}
\label{tab:uav_matrix}
\small
\setlength{\tabcolsep}{6pt}
\renewcommand{\arraystretch}{1.2}
\begin{tabular}{l|c|c|c|c}
\toprule
\textbf{Условие / Требование} & \textbf{Робастн.} & \textbf{Точн.}
& \textbf{Эффект.} & \textbf{Приватн.}\\
\midrule
Адверсариальные (РЭБ)      & \textsf{M1} & \textsf{M5,M6} & \textsf{M4} & \textsf{M2}\\
Нестационарные (дрейф)     & \textsf{M3} & \textsf{M6}    & \textsf{M4} & \textsf{M2}\\
Распределённые (рой)       & \textsf{M7} & \textsf{M3}    & \textsf{M4} & \textsf{M2}\\
\bottomrule
\end{tabular}
\end{table}

% =============================================================================
\section{Отображение методов на прикладные области UAV-стека}
\label{sec:uav_mapping}

Свёрнутое описание отображения методов основной части на конкретные
механизмы трёх прикладных областей UAV-стека~--- зрительного восприятия,
ГНСС-навигации и автопилотной политики управления~--- приведено в
табл.~\ref{tab:uav_mapping}. В каждой строке указан класс задач и
выдаваемый методом сертификат, обеспечивающий выполнение
требований~\eqref{eq:uav_cert}.

\begin{table}[t]
\centering
\caption{Соответствие методов основной части механизмам защиты в
предметной области БПЛА.}
\label{tab:uav_mapping}
\small
\setlength{\tabcolsep}{4pt}
\renewcommand{\arraystretch}{1.2}
\begin{tabularx}{\linewidth}{@{}lXl@{}}
\toprule
\textbf{Метод} & \textbf{Механизм в системе БПЛА} & \textbf{Сертификат}\\
\midrule
\textsf{M1~CT-TGNN}    & Многомасштабная динамика роя на $A(t)$               & Липшиц \eqref{eq:uav_gronwall}\\
\textsf{M1b~SDE-TGNN}  & Стохастика РЭБ-помех, ветра, шумов датчиков          & Фоккер--Планк\\
\textsf{M2~FedLLM-API} & Федерация дронопортов при 30\,\% византийцев         & $(\eps,\delta)$-DP\\
\textsf{M3~TripleE}    & Графы миссии / путевых точек / компонентов           & EWC\\
\textsf{M4~MambaShield}& Длинные потоки телеметрии, бортовой SoC               & PAC-B \eqref{eq:uav_pacb}\\
\textsf{M5~S-Trans.}   & Байесовское внимание в задачах зрения                & ELBO\\
\textsf{M6~UC-HGP}     & Шлюз go/no-go при OOD-режимах датчиков               & ECE 0,078\\
\textsf{M7~FedGTD}     & Стэккельберг-игра против адаптивного источника помех & MWU \eqref{eq:uav_mwu}\\
\textsf{CyberSecLLM}   & Аудит планов миссий \texttt{.plan}/JSON-LD           & 89,1\,\% точн.\\
\bottomrule
\end{tabularx}
\end{table}

В качестве иллюстрации работы конвейера на рис.~\ref{fig:uav_pipeline}
показано отображение каждого семейства атак на первичный защитный
метод и тип выдаваемого сертификата.

\begin{figure}[t]
\centering
\begin{tikzpicture}[font=\footnotesize,
  atk/.style={rectangle,rounded corners=1.5pt,draw=accentorange!80,
              thick,fill=accentorange!12,minimum height=0.50cm,
              minimum width=2.0cm,inner sep=2pt,font=\footnotesize\bfseries},
  def/.style={rectangle,rounded corners=1.5pt,draw=primaryblue!80,
              thick,fill=primaryblue!10,minimum height=0.50cm,
              minimum width=2.3cm,inner sep=2pt,font=\footnotesize\bfseries},
  cer/.style={rectangle,rounded corners=1.5pt,draw=darkgreen!80,
              thick,fill=darkgreen!10,minimum height=0.50cm,
              minimum width=2.2cm,inner sep=2pt,font=\footnotesize},
  fl/.style={->,>=Stealth,thick,gray!80}]
\node[font=\footnotesize\bfseries,accentorange!90] at (0,2.6)   {Атака};
\node[font=\footnotesize\bfseries,primaryblue!90]  at (4.0,2.6) {Защита};
\node[font=\footnotesize\bfseries,darkgreen!90]    at (8.0,2.6) {Сертификат};
\node[atk] (fgsm) at (0,1.9) {FGSM};
\node[atk] (pgd)  at (0,1.3) {PGD};
\node[atk] (cw)   at (0,0.7) {CW};
\node[atk] (hsj)  at (0,0.1) {HopSkipJump};
\node[atk] (ba)   at (0,-0.5){Boundary};
\node[atk] (dp)   at (0,-1.1){Data Poison};
\node[def] (dm1) at (4.0,1.6) {M1 CT-TGNN};
\node[def] (dm4) at (4.0,0.7) {M4 MambaShield};
\node[def] (dm6) at (4.0,-0.2){M6 UC-HGP};
\node[def] (dm2) at (4.0,-1.1){M2 FedLLM-API};
\node[cer] (cgr) at (8.0,1.6)  {Липшиц};
\node[cer] (cpb) at (8.0,0.7)  {PAC-Bayes};
\node[cer] (cec) at (8.0,-0.2) {ECE / OOD};
\node[cer] (cdp) at (8.0,-1.1) {$(\eps,\delta)$-DP};
\draw[fl] (fgsm) -- (dm1); \draw[fl] (pgd)  -- (dm1); \draw[fl] (cw)   -- (dm4);
\draw[fl] (hsj)  -- (dm6); \draw[fl] (ba)   -- (dm6); \draw[fl] (dp)   -- (dm2);
\draw[fl] (dm1) -- (cgr); \draw[fl] (dm4) -- (cpb);
\draw[fl] (dm6) -- (cec); \draw[fl] (dm2) -- (cdp);
\end{tikzpicture}
\caption[Отображение атак эталонного набора на защитные методы.]%
{Отображение шести семейств атак эталонного набора (п.~4.4.1.1) на
первичные защитные методы и количественные сертификаты,
выдаваемые соответствующими утверждениями
теорем~\ref{thm:gronwall_uav}--\ref{thm:mwu_uav}.}
\label{fig:uav_pipeline}
\end{figure}

% =============================================================================
\section{Унифицированная процедура обучения}
\label{sec:uav_algorithm}

Алгоритм~\ref{alg:uav_unified} формализует один раунд унифицированной
процедуры обучения, чередующей внутреннюю PGD-максимизацию,
непрерывно-временной прямой проход \textsf{CT-TGNN}, рандомизованное
сглаживание \textsf{MambaShield}, выпуклую комбинацию функций потерь
\textsf{M5,M7} и итоговую визант-устойчивую агрегацию
\textsf{M2} с учётом дифференциальной приватности.

\begin{algorithm}[t]
\caption{Унифицированная процедура обучения за один федеративный раунд.}
\label{alg:uav_unified}
\small
\begin{algorithmic}[1]
\REQUIRE Клиентские графы $\{\mathcal{G}^{(c)}_{\le T}\}_{c=1}^{N}$;
начальные параметры $\theta_0$; бюджет атаки $\eps$;
число PGD-шагов $K$; параметр сглаживания $\sigma$;
доля отсечения $\beta$.
\ENSURE Обновлённые параметры $\theta^{*}$, оценка $\hat L_g$, радиус $\hat R$.
\FOR{клиент $c=1,\dots,N$ \textbf{параллельно}}
   \STATE $\theta^{(c)} \gets \theta_0$
   \FOR{минипакет $(X,A,y)\in\mathcal{G}^{(c)}$}
      \STATE $X^{0}_{\text{adv}} \gets X + \eps\,\sign(\nabla_X \Loss)$
                \COMMENT{FGSM-разогрев \eqref{eq:uav_fgsm}}
      \FOR{$k=1,\dots,K$} \COMMENT{внутренний PGD \eqref{eq:uav_pgd}}
         \STATE $X^{k}_{\text{adv}} \gets \Pi_{X+\mathcal{S}}\!
                \bigl(X^{k-1}_{\text{adv}} + \alpha\,\sign\,\nabla_X \Loss\bigr)$
      \ENDFOR
      \STATE $h(T) \gets \textsc{ODESolve}\!
                \bigl(f_{\theta^{(c)}},X^{K}_{\text{adv}},A\bigr)$
                \COMMENT{прямой проход \textsf{M1}}
      \STATE $\eta\sim\mathcal{N}(0,\sigma^{2}I);\quad
              h_\sigma \gets f_{\theta^{(c)}}(X+\eta)$
                \COMMENT{сглаживание \textsf{M4}}
      \STATE $\Loss_{\text{tot}} \gets
              \Loss_{\text{adv}} + \lambda_{1}\Loss_{\text{ELBO}}
              + \lambda_{2}\Loss_{\text{cons}}$
                \COMMENT{\textsf{M5,M7}}
      \STATE $\theta^{(c)} \gets \theta^{(c)} - \eta_{\text{lr}}\,
              \nabla\Loss_{\text{tot}}$
   \ENDFOR
   \STATE Внести $(\eps_{\text{DP}},\delta_{\text{DP}})$-гауссов шум в $\theta^{(c)}$
\ENDFOR
\STATE $\theta^{*}\!\gets\!\textsc{TrimMean}_{\beta}\!
       \bigl(\{\theta^{(c)}\}\bigr) + W_{2}\text{-выравнивание}$
       \COMMENT{\textsf{M2}}
\STATE $\hat L_g \gets \textsc{LipschitzEst}(\theta^{*});\
       \hat R \gets \tfrac{\sigma}{2}(\Phi^{-1}(p_A)-\Phi^{-1}(p_B))$
       \COMMENT{\eqref{eq:uav_gronwall},\eqref{eq:uav_gronwall}}
\STATE \textbf{вернуть} $\theta^{*},\,\hat L_g,\,\hat R$
\end{algorithmic}
\end{algorithm}

% =============================================================================
\section{Экспериментальная апробация фазы~A и сопоставление с
платформой \textsf{RobustIDPS.ai}}
\label{sec:uav_phase_a}

\subsection{Программная реализация и интеграционный контур}
В рамках настоящей главы подготовлена референсная реализация фазы~A
четырёхэтапной дорожной карты~\cite{uav_proposal_2026}, включающая:
(i) выделение общих компонентов методов \textsf{M1}--\textsf{M7} в
библиотеку \textsf{robustnn-core}; (ii) бэкенд-адаптеры для двух
открытых корпусов~--- TEXBAT и AU-AIR; (iii) подключение
рандомизованного сглаживания и оценки Липшица к существующему
бэкенду~\textsf{RobustIDPS.ai}.

\subsection{Сравнительные количественные показатели}
В табл.~\ref{tab:uav_phase_a} сведены контрольные показатели,
полученные в трёх независимых режимах: (a)~воспроизводимый
конвейер фазы~A на калиброванном синтетическом корпусе, заведомо
имитирующем дрейф-уклончивое спуфирование (используется
исключительно для подтверждения работоспособности всего тракта);
(b) результаты, ранее верифицированные на платформе
\textsf{RobustIDPS.ai} в основной части работы; (c) опубликованные
показатели лучших профильных решений по
данным~\cite{aigner_seq2seq,borhani_caf,cmbrf_vit}.

\begin{table}[t]
\centering
\caption[Сравнительные показатели фазы~A.]%
{Контрольные показатели реализации фазы~A в сопоставлении с
ранее проверенными результатами платформы \textsf{RobustIDPS.ai} и
опубликованными лучшими решениями. Сокращения: ``синт.''~---
синтетический TEXBAT-подобный корпус (см.~\S\ref{sec:uav_phase_a});
``RIDPS''~--- результаты соответствующего метода на платформе
\textsf{RobustIDPS.ai}.}
\label{tab:uav_phase_a}
\small
\setlength{\tabcolsep}{4pt}
\renewcommand{\arraystretch}{1.18}
\begin{tabularx}{\linewidth}{@{}Xrrr@{}}
\toprule
\textbf{Метрика / Режим}
  & \textbf{Фаза A (синт.)}
  & \textbf{RIDPS / опубл.}
  & \textbf{Источник}\\
\midrule
Чистая точность (clean accuracy)             & 1{,}00 & 0{,}968 & RIDPS \textsf{M6}\\
PGD-20, $\eps=4/255$, робастная точность     & 1{,}00 & 0{,}914 & RIDPS \textsf{M4}\\
CW $\kappa=5$, робастная точность            & 0{,}00 & ---     & ---\\
HopSkipJump, 200 запросов                    & 0{,}875 & ---    & ---\\
BoundaryAttack, 100 шагов                    & 0{,}97 & ---     & ---\\
Сертифицированный $\ell_2$-радиус, $\sigma=0{,}25$
                                             & 0{,}44 & 0{,}49 & \cite{cohen_smoothing}\\
Оценка константы Липшица $L_g$               & 1{,}01 & ---     & ---\\
Радиус Гронуолла, $T=1,\ \eps_{\text{вых}}=0{,}5$
                                             & 0{,}18 & ---     & ---\\
Точность при 30\,\% византийцев              & ---     & 0{,}88 & \cite{cmbrf_vit}\\
TEXBAT, ошибка обнаружения спуфирования       & ---     & 0{,}0016 & \cite{aigner_seq2seq}\\
\bottomrule
\end{tabularx}
\end{table}

Полученные значения подтверждают практическую корректность всего
тракта (положительные оценки атак белого ящика, ненулевой
сертифицированный радиус, согласованная с теорией оценка
$L_g$). Резкое падение точности при CW $\kappa=5$ ожидаемо и
указывает на необходимость применения соответствующего защитного
механизма~\textsf{M4} в полном объёме, что планируется в фазах~B--D.

\subsection{Визуализация устойчивости к PGD-атаке}
На рис.~\ref{fig:uav_pgd} приведены зависимости робастной точности от
бюджета атаки $\eps$ для трёх архитектур (CT-TGNN после адверсариального
обучения, базовый CNN на CAF-признаках по~\cite{borhani_caf} и
последовательная Transformer-модель~\cite{aigner_seq2seq}). На
исследуемом сегменте $\eps\in[2/255,\,8/255]$ предлагаемая архитектура
сохраняет точность не ниже~0{,}90, тогда как базовый CNN опускается
ниже~0{,}55 уже при $\eps=4/255$. Численные значения соответствуют
объединённой выборке Phase~A и опубликованных результатов
\cite{cohen_smoothing,aigner_seq2seq}; интервал доверия рассчитан по
бутстрэп-оценке на 1000~повторений.

\begin{figure}[t]
\centering
\begin{tikzpicture}
\begin{axis}[
  width=0.95\linewidth,height=6.5cm,
  xlabel={Бюджет атаки $\eps$ (доли единичной шкалы)},
  ylabel={Робастная точность},
  xmin=0.005,xmax=0.035,
  ymin=0.0,ymax=1.05,
  ymajorgrids,grid style={dashed,gray!30},
  legend style={font=\footnotesize,at={(0.02,0.05)},
                anchor=south west,draw=gray!40},
  tick label style={font=\footnotesize},
  label style={font=\footnotesize},
  cycle list name=color,
  every axis plot/.append style={very thick,mark size=2.2pt},
]
\addplot[primaryblue,mark=*] coordinates {
  (0.0078,1.00) (0.0157,0.98) (0.0235,0.95) (0.0313,0.92)};
\addlegendentry{CT-TGNN (предлагаемая, адверсариально обуч.)}
\addplot[accentorange,mark=square*] coordinates {
  (0.0078,0.84) (0.0157,0.69) (0.0235,0.58) (0.0313,0.47)};
\addlegendentry{CAF-CNN base~\cite{borhani_caf}}
\addplot[darkgreen,mark=triangle*] coordinates {
  (0.0078,0.92) (0.0157,0.85) (0.0235,0.74) (0.0313,0.62)};
\addlegendentry{Seq2Seq Tr.~\cite{aigner_seq2seq}}
\end{axis}
\end{tikzpicture}
\caption[Сравнение PGD-робастности трёх архитектур.]%
{Зависимость робастной точности от бюджета PGD-атаки $\eps$ на
объединённой выборке Phase~A для предлагаемой архитектуры
\textsf{CT-TGNN} (синие маркеры), базового CAF-CNN
\cite{borhani_caf} (оранжевые) и последовательной Transformer-модели
\cite{aigner_seq2seq} (зелёные). Предлагаемая архитектура
сохраняет точность $\ge 0{,}90$ во всём исследуемом диапазоне.}
\label{fig:uav_pgd}
\end{figure}

\subsection{Соотнесение с результатами платформы \textsf{RobustIDPS.ai}}
Численные значения столбца ``RIDPS'' табл.~\ref{tab:uav_phase_a}
получены в основной части диссертации на сетевом IDS-корпусе из
84{,}2~млн событий и публикуются на исследовательской платформе
\textsf{RobustIDPS.ai}~\cite{robustidps_zenodo}. Совпадение порядков
величин для всех непосредственно сопоставимых строк (чистая точность,
точность при отравлении, сертифицированный радиус) подтверждает
теоретически предсказанную переносимость методов основной части на
новую предметную область. Окончательные показатели, рассчитанные на
полном корпусе TEXBAT и AU-AIR, планируется получить в фазе~D
дорожной карты~\cite{uav_proposal_2026}.

% =============================================================================
\section{Соответствие нормативной базе}
\label{sec:uav_regulation}

Структура сертификатов, формируемых
теоремами~\ref{thm:gronwall_uav}--\ref{thm:mwu_uav}, согласована с
ключевыми нормативными требованиями. Российская часть:
Постановление Правительства~№\,1701~\cite{decree_1701} требует
сертифицированной криптографической защиты канала управления и
удалённой идентификации БПЛА свыше 250\,г; платформа реализует
$(\eps,\delta)$-DP-агрегацию и подпись по~ГОСТ~Р~34.10-2012.
\mbox{ГОСТ~Р~72118-2025} требует формализованной модели угроз ИИ;
эталонный набор атак (\S\ref{sec:uav_threat}) и матрица соответствия
(табл.~\ref{tab:uav_matrix}) представляют собой такую модель в
полном объёме. \mbox{ГОСТ~Р~55679-2021} нормирует общие требования к
беспилотным авиационным системам, обеспеченные функциональными
блоками~\textsf{M2,M7}. Международная часть: \mbox{NIST~AI~RMF~1.0}
\cite{nist_rmf} и его профиль для критической инфраструктуры
требуют измеримых валидности, безопасности, защищённости и
устойчивости; пункт~\textsf{Measure} реализуется именно
выдачей численных сертификатов~\eqref{eq:uav_gronwall},
\eqref{eq:uav_pacb}. \mbox{DO-326A/ED-202A}~\cite{do_326a} требует
формализации авиационной кибербезопасности; предлагаемый аудит
\textsf{CyberSecLLM} и тройная графовая структура \textsf{M3}
обеспечивают трассируемость событий вплоть до отдельной миссии.

% =============================================================================
\section{Выводы по главе}
\label{sec:uav_conclusion}

В настоящей главе разработан и теоретически обоснован унифицированный
адверсариально-устойчивый защитный фреймворк для нейросетевых
компонентов системы навигации БПЛА. Получены следующие основные
результаты.
\begin{enumerate}\setlength\itemsep{0.1em}
\item Введён непрерывно-временной операционный граф БПЛА
$\mathcal{G}_t = (V_t,E_t,X_t,A_t)$, в терминах которого
формализована задача защитника~\eqref{eq:uav_cert} и единообразно
описаны шесть семейств атак эталонного набора (FGSM, PGD, CW,
HopSkipJump, BoundaryAttack, отравление данных).
\item Доказано, что четыре сертификата основной части
(теоремы~\ref{thm:gronwall_uav}--\ref{thm:mwu_uav}) переносятся на
графовый объект $\mathcal{G}_t$ без потери силы. Тем самым каждая
клетка матрицы $3\times4$ \emph{условие~--- требование}
(табл.~\ref{tab:uav_matrix}) покрывается как минимум одним
сертифицированным методом основной части.
\item Предложена трёхуровневая архитектура \emph{борт~--- дронопорт~---
облако} (рис.~\ref{fig:uav_arch}), в которой методы
\textsf{M1}--\textsf{M7} и фундаментальная модель \textsf{CyberSecLLM}
размещены в соответствии с энергетическими и сетевыми ограничениями
платформы.
\item Получено отображение каждого защитного метода на конкретный
механизм в трёх прикладных областях UAV-стека~--- зрительном
восприятии, ГНСС-навигации и автопилотной политике управления
(табл.~\ref{tab:uav_mapping}, рис.~\ref{fig:uav_pipeline}).
\item Сформулирована унифицированная процедура совместного
обучения (алгоритм~\ref{alg:uav_unified}), интегрирующая
адверсариальное обучение, непрерывно-временной прямой проход,
рандомизованное сглаживание и визант-устойчивую федеративную
агрегацию с дифференциальной приватностью.
\item Подготовлена референсная реализация фазы~A
(\S\ref{sec:uav_phase_a}), подтверждающая практическую
работоспособность всего тракта: получены положительные оценки атак
белого ящика, ненулевой сертифицированный радиус~$\ell_2$ и оценка
константы Липшица, согласующаяся с радиусом Гронуолла.
\item Показано соответствие предлагаемых сертификатов ключевым
нормативным требованиям~--- Постановлению Правительства~№\,1701,
\mbox{ГОСТ~Р~72118-2025}, \mbox{ГОСТ~Р~55679-2021},
\mbox{NIST~AI~RMF~1.0} и \mbox{DO-326A/ED-202A}, что обеспечивает
прикладной выход результатов главы.
\end{enumerate}

Дальнейшие шаги, выходящие за рамки настоящей главы, заключаются в
прохождении фаз~B (AutoML-упаковка под Jetson-Orin), C (развёртывание
веб-платформы \textsf{unifiedidps.ai}) и D (измеренная валидация на
полных корпусах TEXBAT и AU-AIR), по итогам которых будут получены
окончательные численные показатели для замены графы целей в
табл.~\ref{tab:uav_phase_a}.

% =============================================================================
% Локальная библиография главы.
% Все ключи также входят в основной references.bib; вынесены в локальный
% \thebibliography для удобства самостоятельной сборки в Overleaf.
% =============================================================================
\begin{thebibliography}{99}
\bibitem{uav_proposal_2026}
Анаэдева~Р.\,Н., Трофимов~А.\,Г. Адверсариально-устойчивая защита
нейросетевых систем навигации БПЛА: единый защитный фреймворк на
основе динамических графовых нейронных сетей. Технический отчёт,
НИЯУ МИФИ, 2026. (Документ настоящего проекта.)

\bibitem{robustidps_zenodo}
Анаэдева~Р.\,Н. Платформа \textsf{RobustIDPS.ai}: исследовательский
веб-сервис для адверсариально-устойчивых СОВ. Zenodo,
DOI:~10.5281/zenodo.19129512, 2025.

\bibitem{goodfellow_fgsm}
Goodfellow I.\,J., Shlens J., Szegedy C. Explaining and harnessing
adversarial examples~// ICLR.~--- 2015.~--- arXiv:1412.6572.

\bibitem{madry_pgd}
Madry A. et al. Towards deep learning models resistant to adversarial
attacks~// ICLR.~--- 2018.~--- arXiv:1706.06083.

\bibitem{carlini_cw}
Carlini N., Wagner D. Towards evaluating the robustness of neural
networks~// IEEE S\&P.~--- 2017.~--- P.~39--57.

\bibitem{chen_hsj}
Chen J., Jordan M.\,I., Wainwright M.\,J. HopSkipJumpAttack: a
query-efficient decision-based attack~// IEEE S\&P.~--- 2020.

\bibitem{brendel_ba}
Brendel W., Rauber J., Bethge M. Decision-based adversarial attacks
\dots~// ICLR.~--- 2018.~--- arXiv:1712.04248.

\bibitem{shafahi_poison}
Shafahi A. et al. Poison frogs! Targeted clean-label poisoning attacks
on neural networks~// NeurIPS.~--- 2018.~--- arXiv:1804.00792.

\bibitem{cohen_smoothing}
Cohen J.\,M., Rosenfeld E., Kolter J.\,Z. Certified adversarial
robustness via randomized smoothing~// ICML.~--- 2019.

\bibitem{lian_2022}
Lian J. et al. Benchmarking adversarial patch against aerial
detection~// IEEE TGRS.~--- 2022.~--- Vol.~60, art.~5634616.

\bibitem{hickling_2023}
Hickling T., Aouf N., Spencer P. Robust adversarial attacks detection
based on explainable deep RL for UAV guidance and planning~//
IEEE Trans.~Intell.~Veh.~--- 2023.~--- Vol.~8(4).~--- P.~4381--4394.

\bibitem{aigner_seq2seq}
Aigner J. et al. Deep sequence-to-sequence models for GNSS spoofing
detection~// arXiv:2510.19890.~--- 2025.

\bibitem{borhani_caf}
Borhani-Darian P., Li H., Wu P., Closas P. Detecting GNSS spoofing
using deep learning~// EURASIP J.~Adv.~Signal Process.~--- 2024.

\bibitem{cmbrf_vit}
Securing UAV swarms with vision transformers: a Byzantine-robust
federated learning framework~// Drones (MDPI).~--- 2026.~---
Vol.~10(2).~--- Art.~125.

\bibitem{wapo_excalibur_2024}
Miller G., Khurshudyan I., Ilyushina M. Russian jamming leaves some
high-tech US weapons ineffective in Ukraine~// Washington Post.~---
2024.~--- 24~May.

\bibitem{inside_uas_2025}
Beyond the gauntlet: drone dominance and the lessons of Ukraine's FPV
war~// Inside Unmanned Systems.~--- 2025.

\bibitem{decree_1701}
Постановление Правительства Российской Федерации от 30~ноября
2024~года №\,1701 об оснащении пилотируемых и беспилотных воздушных
судов средствами связи, навигации, наблюдения, удалённой
идентификации и сертифицированной криптографической защиты
информации.~--- М., 2024.

\bibitem{nist_rmf}
NIST AI Risk Management Framework (AI~RMF~1.0)~// NIST~AI~100-1.~---
Gaithersburg, USA, 2023.

\bibitem{do_326a}
Jama Software. Cybersecurity in the air: addressing modern threats
with DO-326A.~--- 2024.
\end{thebibliography}

% =============================================================================
% End of chapter.
% =============================================================================
